
\medskip
\phantomsection\label{lem:sga3-V-7.2}%
\noindent\textbf{Lemma 7.2.}---%
\emph{Let \(T\) be an affine Noetherian scheme, let \(X,Y,Z\) be
\(T\)-schemes of finite type, with \(X\) quasi-projective over \(T\),
and let}
\SGAThreeDiagram[0.30\textwidth]
  {figures/exp5/p275_lemma72_square.png}
  {The diagram in \hyperref[lem:sga3-V-7.2]{Lemma~7.2}.}
  {fig:sga3-V-7.2-square}
  {Expos\'e V, component-local p.~25 / combined-reader p.~287 / re-edited header p.~275}
\emph{be a diagram in \((\mathbf{Sch}/T)\). Let \(z\) be a point of
\(v(Y)\) that is closed relative to \(T\), and suppose that \(v\) is
flat at the points of \(v^{-1}(z)\). Then there exists a closed
subscheme \(F\) of \(X\) such that
\[
  u\bigl(u^{-1}(F)\cap v^{-1}(z)\bigr)
\]
is finite and nonempty, and such that the restriction of \(v\) to
\(u^{-1}(F)\) is flat at the points of \(v^{-1}(z)\).}

Put \(T=\operatorname{Spec}A\). We may suppose that
\[
  X=\operatorname{Proj}\mathscr S,
\]
where \(\mathscr S\) is the symmetric algebra of a finitely generated
\(A\)-module \(E\).

If \(u(v^{-1}(z))\) is finite, we may take \(F=X\). Otherwise, let
\[
  y_1,\ldots,y_n
\]
be the points of the fiber \(v^{-1}(z)\) associated with its
structure sheaf \(\mathcal O_{v^{-1}(z)}\). Thus, if
\(\mathcal O_i\) denotes the local ring of \(v^{-1}(z)\) at \(y_i\),
the maximal ideal of \(\mathcal O_i\) consists of zero divisors. Let
\(t\) be the image of \(z\) in \(T\). The set \(u(v^{-1}(z))\) is an
infinite constructible subset of the fiber of \(X\) over \(t\).
Consequently, that fiber contains a closed point
\(x\in u(v^{-1}(z))\) distinct from
\(u(y_1),\ldots,u(y_n)\). The open subset
\(X-\{x\}\) is a neighborhood of \(u(y_1),\ldots,u(y_n)\), and hence
contains an open neighborhood of the form \(D_+(f)\), where
\(f\in\mathscr S\) is homogeneous of degree \(d\) (we use the
notation of EGA~II, \S2.3).

The closed subscheme
\[
  X_1=V_+(f)
\]
defined by \(f\) therefore contains \(x\) and avoids the points
\(u(y_1),\ldots,u(y_n)\). Its inverse image
\[
  Y_1=u^{-1}\bigl(V_+(f)\bigr)
\]
is consequently distinct from \(Y\) and meets \(v^{-1}(z)\). We
shall in fact prove that the restriction \(v_1\) of \(v\) to \(Y_1\)
is flat at the points of \(v^{-1}(z)\). If
\(u(v_1^{-1}(z))\) is finite, we may then take \(F=X_1\). Otherwise,
we repeat the argument with \(Y_1,v_1\), and the morphism \(u_1\)
induced by \(u\) on \(Y_1\), in place of \(Y,v,u\). In this way we
obtain a descending sequence
\[
  X\supset X_1\supset\cdots
\]
of closed subschemes of \(X\). Since such a sequence terminates,
\[
  u\bigl(u^{-1}(X_n)\cap v^{-1}(z)\bigr)
\]
is finite and nonempty for some \(n\), and we then take \(F=X_n\).

It remains to show that \(v_1\) is flat at the points of
\(v^{-1}(z)\). Let \(y\in Y_1\) lie over \(z\), let
\(\mathcal O_y\) be the local ring of \(Y\) at \(y\), let
\(\overline{\mathcal O}_y\) be the local ring of \(v^{-1}(z)\) at
\(y\), and let \(\mathcal O_{v(y)}\) be the local ring of \(Z\) at
\(v(y)\). Choose \(g\in\mathscr S_1\) such that \(D_+(g)\) is a
neighborhood of \(u(y)\) in \(X\), and let \(\varphi\) and
\(\overline\varphi\) be the respective images of \(f/g^d\) in
\(\mathcal O_y\) and \(\overline{\mathcal O}_y\). By the construction
of \(f\), the element \(\overline\varphi\) is not a zero divisor in
\(\overline{\mathcal O}_y\). Since \(\mathcal O_y\) is flat over
\(\mathcal O_z=\mathcal O_{v(y)}\), the element \(\varphi\) is not a zero divisor in
\(\mathcal O_y\), and
\[
  \mathcal O_y/\mathcal O_y\varphi
\]
is flat over \(\mathcal O_z\) (SGA~1, IV, 5.7). This quotient is
precisely the local ring of \(Y_1\) at \(y\).
